\documentclass[10pt,a4paper]{article}
\usepackage[margin=2.0cm]{geometry}
\usepackage{amsmath,amssymb,amsthm}
\usepackage[T1]{fontenc}
\usepackage{lmodern}
\usepackage{microtype}
\usepackage[hidelinks]{hyperref}
\usepackage{enumitem}
\setlist{nosep,leftmargin=*}
\theoremstyle{plain}
\newtheorem{theorem}{Theorem}[section]
\newtheorem{proposition}[theorem]{Proposition}
\newtheorem{lemma}[theorem]{Lemma}
\theoremstyle{definition}
\newtheorem{definition}[theorem]{Definition}
\newtheorem{example}[theorem]{Example}
\theoremstyle{remark}
\newtheorem{remark}[theorem]{Remark}
\title{\vspace{-1.2cm}Why Kapustka's degree-20 contraction does not
embed in $\mathbb{P}^7$,\\ and what that does and does not say about
$\mathrm{SR}(\mathcal{M})$}
\author{Note for Daniel \quad (run
\texttt{runs/kapustka-degree20-tangent-space-note-2026-09-18})}
\date{2026-09-18}
\begin{document}
\maketitle
\vspace{-0.9cm}
\noindent\textbf{Scope.} One limited question: the contracted
threefold $\overline{Y}$ in G.~Kapustka's construction of a
degree-20 Calabi--Yau threefold (arXiv:1010.3895, Table~1, No.~8)
has a singular point with a nine-dimensional Zariski tangent
space. Why is that true, why does it stop $\overline{Y}$ from
lying in $\mathbb{P}^7$, and what does it mean for comparing
Kapustka's smooth fibres with a smoothing of the Stanley--Reisner
scheme $\mathrm{SR}(\mathcal{M})$? The smoothing of
$\mathrm{SR}(\mathcal{M})$ is under review and is \emph{not}
assumed here. Every computation quoted below was rerun exactly over
$\mathbb{Q}$ unless it is explicitly marked as a finite-field check;
the logs are in \texttt{logs/}, the sources and quotations in
\texttt{SOURCES.md}.

\section{The objects}

\subsection*{The sphere side}
$\mathcal{M}$ is the Gr\"unbaum--Sreedharan eight-vertex
triangulation of $S^3$. In the labelling of the repository its
Stanley--Reisner ideal in $S=k[a,\dots,h]$ is generated by the
sixteen cubic monomials
$$
\begin{aligned}
&abf,\ abg,\ abh,\ acg,\ ach,\ adh,\ bdf,\ bdg,\\
&beg,\ cde,\ ceg,\ ceh,\ cfh,\ def,\ dfh,\ efg ,
\end{aligned}
$$
and $\mathrm{SR}(\mathcal{M})=\operatorname{Proj} S/I_{\mathcal{M}}
\subset\mathbb{P}^7$ is the union of the twenty coordinate
$\mathbb{P}^3$'s indexed by the facets. Exactly over $\mathbb{Q}$
(\texttt{logs/c01}): $S/I_{\mathcal{M}}$ has Hilbert function
$1,8,36,104,232,440$, Hilbert polynomial
$P(n)=\tfrac{10}{3}n^3+\tfrac{14}{3}n$, degree $20$, no quadrics,
and minimal free resolution with Betti numbers $1,16,30,16,1$ in
degrees $0,3,4,5,8$.

\begin{definition}
A closed subscheme $Z\subset\mathbb{P}^N$ is \emph{arithmetically
Gorenstein} (aG) if its homogeneous coordinate ring $S/I_Z$ is a
Gorenstein ring: it is Cohen--Macaulay (its minimal free resolution
over $S$ has length $\operatorname{codim} Z$) and the last module of
that resolution has rank one. This is a property of the
\emph{embedding}. It is stronger than $Z$ being \emph{Gorenstein},
which is the local property that every local ring of $Z$ is
Gorenstein (smooth points, hypersurface singularities and quotients
of these by groups acting trivially on the canonical form are
Gorenstein).
\end{definition}

So $\mathrm{SR}(\mathcal{M})$ is aG in $\mathbb{P}^7$ with
$h$-vector $(1,4,10,4,1)$, and the last resolution degree $8$ says
$\omega_{\mathrm{SR}(\mathcal{M})}\cong\mathcal{O}$. At a vertex,
say $a=1$, all other seven coordinates are local equations of
degree $\ge 2$ ($\mathcal{M}$ is neighbourly), so the Zariski
tangent space there is all of $\mathbb{A}^7$ (\texttt{logs/c01}):
$\mathrm{SR}(\mathcal{M})$ is as singular as a subscheme of
$\mathbb{P}^7$ can be at those points, but it lives in
$\mathbb{P}^7$ by construction.

\subsection*{Kapustka's construction}
A \emph{del Pezzo surface} is a smooth projective surface with
$-K$ ample; its degree is $K^2$. We need only
$D=\mathbb{P}^1\times\mathbb{P}^1$, with $-K_D=\mathcal{O}(2,2)$,
$K_D^2=8$. The nine bidegree-$(2,2)$ monomials embed $D$ in
$\mathbb{P}^8$ as a surface $D_8$ of degree $8$, cut out by twenty
quadrics (\texttt{logs/c02}); this is the \emph{anticanonical
embedding}. Kapustka projects $D_8$ from a general line of
$\mathbb{P}^8$ to $S\subset\mathbb{P}^6$. The projection is an
isomorphism $D\cong S$ (a general line misses the five-dimensional
secant variety), but $S$ is \emph{not linearly normal}: the seven
coordinates of $\mathbb{P}^6$ restrict to a $7$-dimensional subspace
$V$ of the $9$-dimensional $H^0(S,\mathcal{O}_S(1))=H^0(D,\mathcal{O}(2,2))$.
Write $h^1(I_S(1))=9-7=2$ for the deficit. Over $\mathbb{Q}$, for a
random integer centre (\texttt{logs/c03}): $I_S$ has three quadrics
and fourteen cubics, the Betti table Kapustka prints in the proof of
his Theorem~5.1, and
$$
h^1(I_S(k))=0\quad\text{for } k\ge 2 ,
$$
so $k[S]_n=H^0(D,\mathcal{O}(2n,2n))$ for $n\ge 2$. Since
$h^1(I_S(k))$ is upper semicontinuous in the centre, these values
hold for the general line.

Then: $X'\subset\mathbb{P}^6$ is a general complete intersection of
two quadrics and a cubic containing $S$, a Calabi--Yau threefold of
degree $12$ with $44$ nodes on $S$ (Kapustka, Table~1; the count was
re-derived and finite-field checked in the 2026-09-11 run).
Blowing up $S$ gives a small resolution, and flopping its $44$
exceptional curves gives a smooth Calabi--Yau threefold $X$ with
$\operatorname{Pic}X=\mathbb{Z}H^*\oplus\mathbb{Z}D'$, where $H^*$
pulls back $\mathcal{O}_{X'}(1)$ and $D'\cong D$ is the strict
transform of $S$. Because $K_X=0$, adjunction gives
$$
\mathcal{O}_{D'}(D')=K_{D'},\qquad
\mathcal{O}_{D'}(H^*)=\mathcal{O}_D(2,2)=-K_{D'},
$$
so the divisor $G:=H^*+D'$ restricts trivially to $D'$.
Sections: from $0\to\mathcal{O}_X(H^*)\to\mathcal{O}_X(G)\to
\mathcal{O}_{D'}\to 0$ and $H^1(\mathcal{O}_X(H^*))=0$ (Kawamata--Viehweg,
$H^*$ nef and big, $K_X=0$),
$$
H^0(X,G)=s\cdot H^0(X,H^*)\ \oplus\ \langle w\rangle ,\qquad
\dim=7+1=8,
$$
where $s$ is the section defining $D'$, the seven sections
$s\,x_i$ vanish on $D'$, and $w$ does not. This is the $h^0(H)=8$
of Table~1.

\subsection*{Three different objects}
\begin{enumerate}[label=(\alph*)]
\item The \emph{normal contraction} $\varphi:X\to\overline{Y}$: the
morphism with connected fibres that contracts exactly $D'$ to a
point $P$ (Kapustka takes $\varphi=\varphi_{|2G|}$ in Theorem~3.4
and says in the proof of Theorem~5.1 that the degree-20 case
``concludes as in the proof of Theorem~3.4''). It is a
\emph{primitive contraction of type II} in the sense of Wilson and
Gross: birational onto a normal variety, not factorisable, with
exceptional locus a divisor mapping to a point. There is an ample
Cartier divisor $T$ on $\overline{Y}$ with $\varphi^*T=G$, so
$H^0(\overline{Y},T)=H^0(X,G)$ has dimension $8$
(Kapustka, proof of Theorem~3.4: ``the image $T$ of $G$ on $Y$
is an ample divisor ($2T$ is very ample)'').
\item The \emph{image under the eight sections}
$Y:=\varphi_{|T|}(\overline{Y})\subset\mathbb{P}^7$.
\item The \emph{smooth fibres} $Y_t$ of a smoothing
$\mathcal{Y}\to\Delta$ of $\overline{Y}$, which exists by Gross,
\emph{Deforming Calabi--Yau threefolds}, Theorem~5.8: a primitive
type II contraction is smoothable unless the exceptional divisor is
$\mathbb{P}^2$ or $\mathbb{F}_1$. Table~1 records
$H^3=20$, $h^0(H)=8$, $\chi=-60$, and $h^{1,1}=1$, $h^{1,2}=31$.
\end{enumerate}
Kapustka's Remark~3.3, for the analogous degree-15 row, says
exactly that (b) is not normal at $P$ and that one must pass to a
multiple of $G$ to obtain (a). The point of this note is what
happens at $P$.

\section{The tangent space at $P$ is nine-dimensional}

Recall (HP, algebraic geometry, ``Zariski tangent space'') that for
a point $P$ with local ring $(\mathcal{O}_P,\mathfrak{m})$ the
Zariski tangent space is $T_P=(\mathfrak{m}/\mathfrak{m}^2)^\vee$.
Its dimension is the \emph{embedding dimension}: if $P\in Z$ and
$Z$ is a closed subscheme of a smooth $N$-fold near $P$, then the
cotangent space of the ambient space surjects onto
$\mathfrak{m}/\mathfrak{m}^2$, so $\dim T_PZ\le N$.

\begin{theorem}\label{thm:tangent}
$\dim T_P\overline{Y}=9$ and $\operatorname{mult}_P\overline{Y}=8$.
More precisely, the germ $(\overline{Y},P)$ is analytically the
affine cone over $D_8\subset\mathbb{P}^8$, that is,
$$
(\overline{Y},P)\ \cong\ \bigl(\{xy-zw=0\}\subset\mathbb{A}^4\bigr)/\{\pm 1\},
$$
the three-dimensional ordinary double point divided by the sign
involution.
\end{theorem}

\noindent\textbf{Step 1: the lower bound $\dim T_P\overline{Y}\ge 9$
by hand.}
Let $\widehat{\mathcal{O}}$ be the completion of
$\mathcal{O}_{\overline{Y},P}$. The theorem on formal functions
(Hartshorne III.11.1), applied to the proper morphism $\varphi$
whose fibre over $P$ is $D'$, gives
$$
\widehat{\mathcal{O}}=\varprojlim_n
H^0\bigl(X,\ \mathcal{O}_X/\mathcal{O}_X(-nD')\bigr).
$$
(The theorem uses the powers of $\mathfrak{m}\mathcal{O}_X$; these
are cofinal with the powers of the ideal sheaf $\mathcal{O}_X(-D')$
of the reduced fibre, so the limits agree.) For $n\ge 1$ the
quotient $\mathcal{O}_X(-nD')/\mathcal{O}_X(-(n+1)D')$ is the line
bundle $\mathcal{O}_{D'}(-nD')=\mathcal{O}_D(2n,2n)$, whose $H^1$
vanishes; so all transition maps in the inverse system are
surjective and so is $\widehat{\mathcal{O}}\to
H^0(\mathcal{O}_X/\mathcal{O}_X(-2D'))$. Put
$F_n:=\ker\bigl(\widehat{\mathcal{O}}\to
H^0(\mathcal{O}_X/\mathcal{O}_X(-nD'))\bigr)$, the functions
vanishing to order $\ge n$ along $D'$. Then $F_1=\widehat{\mathfrak{m}}$
(the kernel of evaluation at $P$), $F_1F_1\subseteq F_2$, and
$$
F_1/F_2\ \cong\ H^0\bigl(D',\mathcal{O}_{D'}(-D')\bigr)
= H^0(D,\mathcal{O}(2,2)),\qquad \dim = 9 .
$$
Since $\widehat{\mathfrak{m}}^2\subseteq F_2$, the cotangent space
$\widehat{\mathfrak{m}}/\widehat{\mathfrak{m}}^2$ surjects onto
$F_1/F_2$, whence $\dim T_P\overline{Y}\ge 9$. This inequality is
all that the $\mathbb{P}^7$ obstruction needs.

\noindent\textbf{Step 2: equality, and the published theorem that
identifies the germ.} The same $H^1$-vanishing gives
$F_n/F_{n+1}\cong H^0(D,\mathcal{O}(2n,2n))$ for all $n$, with the
multiplication of sections, so the graded ring
$\bigoplus_n F_n/F_{n+1}$ is the anticanonical ring
$A=\bigoplus_n H^0(D,\mathcal{O}(2n,2n))$ of $D$. It is generated
in degree one (monomials of bidegree $(2n,2n)$ are products of
monomials of bidegree $(2,2)$; \texttt{logs/c02} records the
Hilbert function $(2n+1)^2$). Hence
$F_n=\widehat{\mathfrak{m}}^n+F_{n+1}$ for all $n$, and since the
$F_n$ define the $\widehat{\mathfrak{m}}$-adic topology (Artin--Rees,
as in the proof of the theorem on formal functions), $F_2=\widehat{\mathfrak{m}}^2$:
$\dim T_P=9$ and $\operatorname{mult}_P\overline{Y}=\deg D_8=8$.
The cleanest published route is Gross, \emph{loc.\ cit.},
Proposition~5.4: \emph{an isolated rational Gorenstein threefold
point of multiplicity $k\ge 5$ whose blow-up is smooth with smooth
exceptional divisor $E$ is analytically the cone over $E$}; its
proof uses Grauert's linearisation theorem and
$I_E/I_E^2=\omega_E^{-1}$, so ``cone over $E$'' means the
anticanonical cone. The hypotheses hold for $(\overline{Y},P)$: it
is Gorenstein with $K_{\overline{Y}}=0$ (a Calabi--Yau threefold in
Kapustka's sense), hence rational, and by Reid's theorem quoted in
the proof of Gross's Theorem~5.2 a primitive type II contraction is
the blow-up of the singular point, here with $\widetilde{X}=X$ and
$E=D'\cong\mathbb{P}^1\times\mathbb{P}^1$ of degree $8$. This is
the theorem that connects the local calculation to Kapustka's
contraction.

\noindent\textbf{Step 3: the cone is the quotient of a node.}
Write $D=\{q=0\}\subset\mathbb{P}^3=\mathbb{P}(U)$ with
$q=xy-zw$; then $-K_D=\mathcal{O}_{\mathbb{P}^3}(2)|_D$, so
$A_n=H^0(D,\mathcal{O}(2n))$ is the even part of $k[x,y,z,w]/(q)$,
that is, the ring of invariants of $\pm 1$ acting on the node
$\{q=0\}\subset\mathbb{A}^4$. Exactly over $\mathbb{Q}$
(\texttt{logs/c02}): the kernel of $k[a_{ij}]\to k[x,y,z,w]/(q)$,
$a_{ij}\mapsto$ (product of the $i$th and $j$th coordinates), is the
one linear form $\lambda=a_{01}-a_{23}$ together with twenty
quadrics, and modulo $\lambda$ it coincides with the ideal of $D_8$
after the parametrisation $(x,y,z,w)=(su,tv,sv,tu)$. The vertex has
tangent dimension $10-1=9$; every other point of the cone has
tangent dimension $3$. \hfill$\square$

\section{Consequences, and why they do not survive smoothing}

\begin{proposition}\label{prop:noembed}
$\overline{Y}$ has no closed embedding into $\mathbb{P}^7$, or into
any smooth $7$-fold. In particular $T$ is not very ample. The
morphism $\varphi_{|T|}:\overline{Y}\to\mathbb{P}^7$ defined by the
eight sections has differential of rank $\le 7$ at $P$, and its
image $Y$ is not normal at $\varphi_{|T|}(P)$: the normalisation
$\nu:\overline{Y}\to Y$ satisfies
$\dim_k\nu_*\mathcal{O}_{\overline{Y}}/\mathcal{O}_Y=2$, so
$$
\chi(\mathcal{O}_Y(n))=P(n)-2=\tfrac{10}{3}n^3+\tfrac{14}{3}n-2 .
$$
Hence $[Y]$ and $[\mathrm{SR}(\mathcal{M})]$ lie in different
Hilbert schemes of $\mathbb{P}^7$.
\end{proposition}

\begin{proof}
The first sentence is $\dim T_P\overline{Y}=9>7$ together with the
surjection of cotangent spaces recalled above. ``Very ample'' means
that the complete linear system $|T|$, here of dimension $7$,
embeds $\overline{Y}$; so $T$ is not very ample. Eight sections
give a morphism to $\mathbb{P}^7$ because $|T|$ is base point free:
the seven sections $s\,x_i$ vanish exactly on $D'$ (the $x_i$ have
no common zero on $X'$) and $w$ does not vanish on $D'$. At $P$,
in the chart $w\ne 0$, the seven functions $x_i s/w$ lie in
$\widehat{\mathfrak{m}}$ and their images in
$F_1/F_2=H^0(D,\mathcal{O}(2,2))$ are the seven linear forms of
$\mathbb{P}^6$ restricted to $S$, spanning the $7$-dimensional
$V$. So the differential misses a $2$-dimensional part of $T_P$:
this is precisely why eight sections do not embed. For the
non-normality, $\widehat{\mathcal{O}}_Y$ is the closed subalgebra
of $\widehat{\mathcal{O}}$ generated by those seven functions, and
with the filtration induced from $F_\bullet$ its graded ring is
$k\oplus V\oplus A_2\oplus A_3\oplus\cdots$: in degree one only $V$,
in degree $n\ge 2$ all of $A_n$ because $k[S]_n=A_n$, that is,
$h^1(I_S(n))=0$ for $n\ge 2$ (\texttt{logs/c03}). Hence
$\dim\widehat{\mathcal{O}}/\widehat{\mathcal{O}}_Y=\dim A_1/V=2$.
Finally $\chi(\mathcal{O}_{\overline{Y}}(nT))$ is a cubic
polynomial with leading coefficient $T^3/6=20/6$, constant term
$\chi(\mathcal{O}_{\overline{Y}})=0$, odd in $n$ by Serre duality
($\omega=\mathcal{O}$), and value $8$ at $n=1$ (vanishing of higher
cohomology of $T$ on $\overline{Y}$, which has rational
singularities); so it equals $P(n)$, and
$\chi(\mathcal{O}_Y(n))=P(n)-2$ because
$\nu_*\mathcal{O}_{\overline{Y}}/\mathcal{O}_Y$ is a skyscraper of
length $2$. The 2026-09-11 run's finite-field model of $Y$
(\texttt{logs/c04}, over $\mathbb{F}_{32003}$) confirms
$\dim T_PY=7$ and the Hilbert polynomial $P(n)-2$.
\end{proof}

The same fact seen globally: the section ring
$R=\bigoplus_n H^0(\overline{Y},nT)$ has the Hilbert function
$1,8,36,104,\dots$ of $k[\mathcal{M}]$, but needs two generators
in degree $2$ besides the eight in degree $1$, because
$\operatorname{Sym}^2H^0(T)\to H^0(2T)$, a map of $36$-dimensional
spaces, kills the two quadrics of $X'$. Locally, two of the nine
cotangent directions at $P$ are invisible to the eight sections.
Both are the two units of linear normality lost when $D_8$ was
projected from a line. $R$ is a Gorenstein ring (vanishing plus
$K_{\overline{Y}}=0$), so $\overline{Y}$ is ``arithmetically
Gorenstein'' in the weighted space $\operatorname{Proj}R$, but not
in $\mathbb{P}^7$, where its image is not even normal.

\subsection*{Why none of this persists after smoothing}
The function $t\mapsto\max_{y\in Y_t}\dim T_yY_t$ is upper
semicontinuous: tangent dimensions can only drop when one moves
away from the special fibre, and on a smooth fibre they equal $3$
everywhere.

\begin{example}
The surface node $\{xy=z^2\}\subset\mathbb{A}^3$ has
$\dim T_0=3$; the fibres $\{xy-z^2=t\}$, $t\ne 0$, are smooth with
$\dim T=2$ at every point. Closer to home, Gross (\emph{loc.\
cit.}~5.5) smooths the cone of Theorem~\ref{thm:tangent} by
hyperplane sections of the cone over the second Veronese
$v_2(\mathbb{P}^3)\subset\mathbb{P}^9$. Let $W_t\subset\mathbb{A}^{10}$
be the set of symmetric $4\times 4$ matrices $A$ of rank $\le 1$
with $\lambda(A)=t$, where $\lambda(A)=a_{01}-a_{23}$. Then
$W_0=(\overline{Y},P)$, and for $t\ne 0$, $W_t=\{xy-zw=t\}/\{\pm1\}$
is smooth because $\pm1$ acts freely on the smooth affine quadric.
Over $\mathbb{Q}$ (\texttt{logs/c02}): $\dim T=9$ at the vertex of
$W_0$ and $\dim T=3$ at two sample points of $W_1$. Embedding
dimension $9$ becomes $3$ the moment $t\ne 0$.
\end{example}

The property that matters for $Y_t$ is therefore not local. It is
whether the \emph{eight} sections of $T_t$ embed $Y_t$ with
projectively normal image, that is, whether
$$
(\star)\qquad
\mu_t:\operatorname{Sym}^2H^0(Y_t,T_t)\longrightarrow H^0(Y_t,2T_t)
\quad(36\to 36)
$$
is an isomorphism, equivalently whether $Y_t$ is aG in
$\mathbb{P}^7$. Very ampleness and projective normality are open
conditions in the family: they fail at $t=0$ (Proposition~\ref{prop:noembed}),
and openness gives no information in the other direction.
Semicontinuity gives $\dim\ker\mu_t\le\dim\ker\mu_0=2$, which is
the wrong inequality. The singular model is silent on $(\star)$.

\section{The limited relevance to $\mathrm{SR}(\mathcal{M})$}

\begin{proposition}\label{prop:sr}
Let $\mathcal{X}\to(B,0)$ be a flat family of closed subschemes of
$\mathbb{P}^7$ with $\mathcal{X}_0=\mathrm{SR}(\mathcal{M})$ (an
embedded deformation). Then every fibre $\mathcal{X}_b$, $b$ near
$0$, is aG with $h$-vector $(1,4,10,4,1)$, lies on no quadric, is
cut out by $16$ cubics, and has
$\omega_{\mathcal{X}_b}\cong\mathcal{O}_{\mathcal{X}_b}$. If it is
smooth and connected it is a Calabi--Yau threefold of degree $20$
with $h^0(\mathcal{O}(1))=8$, $H^3=20$, $c_2\cdot H=56$.
\end{proposition}

\begin{proof}
Flatness fixes the Hilbert polynomial $P$. Write $H_b(k)$ for the
Hilbert function of $S/I_{\mathcal{X}_b}$. Semicontinuity gives
$\dim(I_{\mathcal{X}_b})_k\le\dim(I_{\mathcal{M}})_k$ and
$h^0(\mathcal{O}_{\mathcal{X}_b}(k))\le h^0(\mathcal{O}_{\mathcal{X}_0}(k))=H_0(k)$
(the last equality because $k[\mathcal{M}]$ is Cohen--Macaulay), so
$H_0(k)\le H_b(k)\le h^0(\mathcal{O}_{\mathcal{X}_b}(k))\le H_0(k)$:
the Hilbert function is constant and the fibres are projectively
normal. With constant Hilbert function the graded Betti numbers
are upper semicontinuous,
$\beta_{i,j}(\mathcal{X}_b)\le\beta_{i,j}(\mathrm{SR}(\mathcal{M}))$,
so the resolution of $\mathcal{X}_b$ has length $\le 4=\operatorname{codim}$
(hence exactly $4$, Cohen--Macaulay) and last Betti number $\le 1$
(hence $1$: Gorenstein), sitting in degree $8$. This gives the
$h$-vector, the absence of quadrics, the $16$ cubics, and
$\omega\cong\mathcal{O}$.
The invariants follow from the Hilbert polynomial by Riemann--Roch
on a smooth fibre.
\end{proof}

The 2026-09-11 run's \texttt{REPORT.md}, Theorem~5.1, asserts more:
that every \emph{abstract} smoothing of $\mathrm{SR}(\mathcal{M})$
is of this embedded form, because $h^1(\mathcal{O})=h^2(\mathcal{O})=0$
(Hochster's formula, $\mathcal{M}$ a homology sphere) lets
$\mathcal{O}(1)$ extend uniquely to the family. I have not
re-verified that step; it is not needed for the smoothing under
review, which is constructed as a deformation of the ideal
$I_{\mathcal{M}}$ inside $\mathbb{P}^7$.

\medskip
What Theorem~\ref{thm:tangent} and Proposition~\ref{prop:noembed}
actually give for the comparison:

\begin{enumerate}[label=(\roman*)]
\item The two \emph{singular} models are incomparable. $\overline{Y}$
is irreducible and normal with one point of embedding dimension
$9$ and is not a subscheme of $\mathbb{P}^7$; its $\mathbb{P}^7$
image $Y$ has Hilbert polynomial $P(n)-2$. $\mathrm{SR}(\mathcal{M})$
is a reducible aG subscheme of $\mathbb{P}^7$ with Hilbert
polynomial $P(n)$. No flat family inside $\mathbb{P}^7$ connects
$Y$ or $\overline{Y}$ to $\mathrm{SR}(\mathcal{M})$, so no
degeneration of Kapustka's construction to
$\mathrm{SR}(\mathcal{M})$ can be read off from the objects
Kapustka writes down. (This is a structural fact, not a failed
search.)
\item Any comparison must go through the smooth fibres, and there it
is exactly $(\star)$. \emph{Conditionally}: if $(\star)$ holds for
general $t$, then $Y_t\subset\mathbb{P}^7$ is an aG Calabi--Yau
threefold of degree $20$ with $h$-vector $(1,4,10,4,1)$ and
$h^{1,1}=1$, hence a point of the same Hilbert scheme as any
embedded smoothing of $\mathrm{SR}(\mathcal{M})$, and the question
becomes whether the two lie on a common component. If $(\star)$
fails, no embedded smoothing of $\mathrm{SR}(\mathcal{M})$ has
fibres isomorphic to Kapustka's $Y_t$ \emph{with}
$T_t=\mathcal{O}(1)$, by Proposition~\ref{prop:sr}. Neither
alternative is decided here, and nothing here says that
Kapustka's $Y_t$ is non-aG or abstractly deformation-inequivalent
to a fibre of an $\mathrm{SR}(\mathcal{M})$ smoothing.
\item The one published aG degree-20 family in $\mathbb{P}^7$ (the
determinantal Gulliksen--Neg\aa rd threefold, HP ``GN variety'') has
$h^{1,1}=2$, so it is a different family from Kapustka's
$h^{1,1}=1$ one, whatever the answer to $(\star)$.
\end{enumerate}

\section*{Conclusion and next question}

\noindent\textbf{Conclusion.} Kapustka's normal contraction
$\overline{Y}$ has a single singular point $P$, analytically the
cone over the anticanonical $\mathbb{P}^1\times\mathbb{P}^1$, i.e.\
$\{xy=zw\}/\{\pm1\}$, with $\dim T_P\overline{Y}=9$ (the lower bound
$\ge 9$ is elementary from the theorem on formal functions and
$H^1(\mathcal{O}_D(2,2))=0$; equality and the analytic model are
Gross's Proposition~5.4). Hence $\overline{Y}$ embeds in no smooth
$7$-fold, $T$ is not very ample, the eight sections map
$\overline{Y}$ onto a non-normal $Y\subset\mathbb{P}^7$ with
Hilbert polynomial $P(n)-2$, and the singular models of the two
constructions cannot be compared inside $\mathbb{P}^7$. This is a
statement about the special fibre only; embedding dimension drops
to $3$ on every smooth fibre, and whether the eight sections of
$T_t$ embed $Y_t$ as an aG threefold in $\mathbb{P}^7$ (property
$(\star)$) is neither implied nor excluded by it.

\noindent\textbf{Next question, for Kapustka.} In Theorem~3.4 the
degree-15 fibres are embedded by $|2T_t|$. For row 8, is
$\mu_t:\operatorname{Sym}^2H^0(T_t)\to H^0(2T_t)$ an isomorphism
for general $t$, i.e.\ does $|T_t|$ embed $Y_t$ in $\mathbb{P}^7$
with projectively normal image? Equivalently: does the general
member of the family lie on no quadric of $\mathbb{P}^7$? The
2026-09-11 run claims a first-order obstruction inside the weighted
model (not re-verified here) and reduces the decision to a
second-order term.

\end{document}
